%───────────────────────────────────────────────────────────────────────────────
\subsection{Adaptive Weighting}
\label{sec:adaptive}

Each loss component is independently normalised by its own magnitude to prevent
either head from dominating the gradient.  For a generic per-sample raw loss
$\ell^{(b)} \in \mathbb{R}_{\ge 0}$, the adaptive weight and weighted loss
are:
\begin{equation}
  w^{(b)}
    \coloneqq \bigl(\sg[\ell^{(b)}] + \varepsilon_{\mathrm{norm}}\bigr)^{p_{\mathrm{norm}}}
    \;\in \mathbb{R}_{> 0}
  \label{eq:adaptive_weight}
\end{equation}
\begin{equation}
  \ell^{(b)}_{\mathrm{weighted}}
    \coloneqq \frac{\ell^{(b)}}{w^{(b)}}
    \;\in \mathbb{R}_{\ge 0}
  \label{eq:adaptive_weighted_loss}
\end{equation}
where $\sg[\cdot]$ is the stop-gradient operator (Eq.~\ref{eq:sg_def})
ensuring the weight does not contribute gradients, $\varepsilon_{\mathrm{norm}}
= 1.0$, and $p_{\mathrm{norm}} = 1.0$.  Substituting these defaults:
\begin{equation}
  w^{(b)} = \sg[\ell^{(b)}] + 1.0,
  \qquad
  \ell^{(b)}_{\mathrm{weighted}}
    = \frac{\ell^{(b)}}{\sg[\ell^{(b)}] + 1.0}
  \label{eq:adaptive_simplified}
\end{equation}
\begin{shapes}
$\ell^{(b)}$: per-sample scalar; full batch: $(B,)$.\\
$w^{(b)}$: per-sample scalar; full batch: $(B,)$. Computed with \code{.detach()} on $\ell$.\\
$\ell^{(b)}_{\mathrm{weighted}}$: per-sample scalar; full batch: $(B,)$.\\
\textbf{Code:} \code{meanflow\_loss.py:347--351}:
\code{adp\_weight\_u = (raw\_loss\_u.detach() + norm\_eps) ** norm\_p};
\code{loss\_u = raw\_loss\_u / adp\_weight\_u}.
Applied \textbf{independently} to u-head and v-head.
\end{shapes}

\noindent
This is a form of \textbf{logarithmic loss normalisation}: the effective loss
behaves as $\ell_\mathrm{w} \approx \ell/(\ell+1)$, which is bounded in
$[0,1)$ and varies slowly with raw loss magnitude:
\begin{equation}
  \lim_{\ell \to 0}\;\frac{\ell}{\ell + 1} = 0,
  \qquad
  \lim_{\ell \to \infty}\;\frac{\ell}{\ell + 1} = 1
\end{equation}

At convergence (epoch~388, best FID), the loss decomposition shows this
balancing in action (Table~\ref{tab:loss_decomp}).

\begin{table}[ht]
\centering
\caption{Loss decomposition at best FID (epoch~388).}
\label{tab:loss_decomp}
\begin{tabular}{lrr}
\toprule
\textbf{Component} & \textbf{Raw Loss} & \textbf{Share} \\
\midrule
FM loss ($r=t$, flow matching)       & 715,432   & 10.3\% \\
MF loss ($r<t$, self-consistency)    & 6,254,997 & 89.7\% \\
\midrule
v-head loss                          & 714,646   & --- \\
u-head loss (compound $V$)           & 3,485,215 & --- \\
\bottomrule
\end{tabular}
\end{table}

The MF loss is ${\sim}9\times$ larger than the FM loss, which is expected: the
compound velocity $V$ must also capture the JVP correction term, a harder
target than the direct velocity at $r{=}t$.  The adaptive weighting equalises
their gradient contributions to ${\sim}1.0$ each.

\textbf{Why $\varepsilon_{\mathrm{norm}}{=}1.0$?}  Through Phase~4c debugging,
we discovered that $\varepsilon_{\mathrm{norm}}{=}0.01$ caused catastrophic
gradient amplification for samples with small raw loss.  The weight
$1/(\ell+0.01)$ can reach $100\times$, creating a positive feedback loop.

\textbf{Why $p_{\mathrm{norm}}{=}1.0$?}  With $p_{\mathrm{norm}}{=}0.5$, the
weight becomes $w = \sqrt{\ell+1.0}$, which under-normalises large
losses, leading to a $1000\times$ gradient explosion (Phase~4e).
